% Options for packages loaded elsewhere
\PassOptionsToPackage{unicode}{hyperref}
\PassOptionsToPackage{hyphens}{url}
\documentclass[
  ignorenonframetext,
]{beamer}
\newif\ifbibliography
\usepackage{pgfpages}
\setbeamertemplate{caption}[numbered]
\setbeamertemplate{caption label separator}{: }
\setbeamercolor{caption name}{fg=normal text.fg}
\beamertemplatenavigationsymbolsempty
% remove section numbering
\setbeamertemplate{part page}{
  \centering
  \begin{beamercolorbox}[sep=16pt,center]{part title}
    \usebeamerfont{part title}\insertpart\par
  \end{beamercolorbox}
}
\setbeamertemplate{section page}{
  \centering
  \begin{beamercolorbox}[sep=12pt,center]{section title}
    \usebeamerfont{section title}\insertsection\par
  \end{beamercolorbox}
}
\setbeamertemplate{subsection page}{
  \centering
  \begin{beamercolorbox}[sep=8pt,center]{subsection title}
    \usebeamerfont{subsection title}\insertsubsection\par
  \end{beamercolorbox}
}
% Prevent slide breaks in the middle of a paragraph
\widowpenalties 1 10000
\raggedbottom
\AtBeginPart{
  \frame{\partpage}
}
\AtBeginSection{
  \ifbibliography
  \else
    \frame{\sectionpage}
  \fi
}
\AtBeginSubsection{
  \frame{\subsectionpage}
}
\usepackage{iftex}
\ifPDFTeX
  \usepackage[T1]{fontenc}
  \usepackage[utf8]{inputenc}
  \usepackage{textcomp} % provide euro and other symbols
\else % if luatex or xetex
  \usepackage{unicode-math} % this also loads fontspec
  \defaultfontfeatures{Scale=MatchLowercase}
  \defaultfontfeatures[\rmfamily]{Ligatures=TeX,Scale=1}
\fi
\usepackage{lmodern}
\usetheme[]{Montpellier}
\ifPDFTeX\else
  % xetex/luatex font selection
\fi
% Use upquote if available, for straight quotes in verbatim environments
\IfFileExists{upquote.sty}{\usepackage{upquote}}{}
\IfFileExists{microtype.sty}{% use microtype if available
  \usepackage[]{microtype}
  \UseMicrotypeSet[protrusion]{basicmath} % disable protrusion for tt fonts
}{}
\makeatletter
\@ifundefined{KOMAClassName}{% if non-KOMA class
  \IfFileExists{parskip.sty}{%
    \usepackage{parskip}
  }{% else
    \setlength{\parindent}{0pt}
    \setlength{\parskip}{6pt plus 2pt minus 1pt}}
}{% if KOMA class
  \KOMAoptions{parskip=half}}
\makeatother
\usepackage{color}
\usepackage{fancyvrb}
\newcommand{\VerbBar}{|}
\newcommand{\VERB}{\Verb[commandchars=\\\{\}]}
\DefineVerbatimEnvironment{Highlighting}{Verbatim}{commandchars=\\\{\}}
% Add ',fontsize=\small' for more characters per line
\usepackage{framed}
\definecolor{shadecolor}{RGB}{248,248,248}
\newenvironment{Shaded}{\begin{snugshade}}{\end{snugshade}}
\newcommand{\AlertTok}[1]{\textcolor[rgb]{0.94,0.16,0.16}{#1}}
\newcommand{\AnnotationTok}[1]{\textcolor[rgb]{0.56,0.35,0.01}{\textbf{\textit{#1}}}}
\newcommand{\AttributeTok}[1]{\textcolor[rgb]{0.13,0.29,0.53}{#1}}
\newcommand{\BaseNTok}[1]{\textcolor[rgb]{0.00,0.00,0.81}{#1}}
\newcommand{\BuiltInTok}[1]{#1}
\newcommand{\CharTok}[1]{\textcolor[rgb]{0.31,0.60,0.02}{#1}}
\newcommand{\CommentTok}[1]{\textcolor[rgb]{0.56,0.35,0.01}{\textit{#1}}}
\newcommand{\CommentVarTok}[1]{\textcolor[rgb]{0.56,0.35,0.01}{\textbf{\textit{#1}}}}
\newcommand{\ConstantTok}[1]{\textcolor[rgb]{0.56,0.35,0.01}{#1}}
\newcommand{\ControlFlowTok}[1]{\textcolor[rgb]{0.13,0.29,0.53}{\textbf{#1}}}
\newcommand{\DataTypeTok}[1]{\textcolor[rgb]{0.13,0.29,0.53}{#1}}
\newcommand{\DecValTok}[1]{\textcolor[rgb]{0.00,0.00,0.81}{#1}}
\newcommand{\DocumentationTok}[1]{\textcolor[rgb]{0.56,0.35,0.01}{\textbf{\textit{#1}}}}
\newcommand{\ErrorTok}[1]{\textcolor[rgb]{0.64,0.00,0.00}{\textbf{#1}}}
\newcommand{\ExtensionTok}[1]{#1}
\newcommand{\FloatTok}[1]{\textcolor[rgb]{0.00,0.00,0.81}{#1}}
\newcommand{\FunctionTok}[1]{\textcolor[rgb]{0.13,0.29,0.53}{\textbf{#1}}}
\newcommand{\ImportTok}[1]{#1}
\newcommand{\InformationTok}[1]{\textcolor[rgb]{0.56,0.35,0.01}{\textbf{\textit{#1}}}}
\newcommand{\KeywordTok}[1]{\textcolor[rgb]{0.13,0.29,0.53}{\textbf{#1}}}
\newcommand{\NormalTok}[1]{#1}
\newcommand{\OperatorTok}[1]{\textcolor[rgb]{0.81,0.36,0.00}{\textbf{#1}}}
\newcommand{\OtherTok}[1]{\textcolor[rgb]{0.56,0.35,0.01}{#1}}
\newcommand{\PreprocessorTok}[1]{\textcolor[rgb]{0.56,0.35,0.01}{\textit{#1}}}
\newcommand{\RegionMarkerTok}[1]{#1}
\newcommand{\SpecialCharTok}[1]{\textcolor[rgb]{0.81,0.36,0.00}{\textbf{#1}}}
\newcommand{\SpecialStringTok}[1]{\textcolor[rgb]{0.31,0.60,0.02}{#1}}
\newcommand{\StringTok}[1]{\textcolor[rgb]{0.31,0.60,0.02}{#1}}
\newcommand{\VariableTok}[1]{\textcolor[rgb]{0.00,0.00,0.00}{#1}}
\newcommand{\VerbatimStringTok}[1]{\textcolor[rgb]{0.31,0.60,0.02}{#1}}
\newcommand{\WarningTok}[1]{\textcolor[rgb]{0.56,0.35,0.01}{\textbf{\textit{#1}}}}
\usepackage{longtable,booktabs,array}
\usepackage{calc} % for calculating minipage widths
\usepackage{caption}
% Make caption package work with longtable
\makeatletter
\def\fnum@table{\tablename~\thetable}
\makeatother
\setlength{\emergencystretch}{3em} % prevent overfull lines
\providecommand{\tightlist}{%
  \setlength{\itemsep}{0pt}\setlength{\parskip}{0pt}}
%\documentclass[unicode,10pt]{beamer}% 'unicode'が必要
\usepackage{luatexja}% 日本語を使う場合は必須
%\usepackage[japanese]{babel}	
%\usefonttheme[onlymath]{serif}
%\usepackage[T1]{fontenc} %これを使うと、ギリシャ文字が出ない
%\usepackage{textcomp}
%\usepackage[scale=1.0]{tgheros} % Sans serif font
%\usepackage[scaled]{beramono} % Monospace font
%\usepackage{luatexja-otf}
%\renewcommand{\kanjifamilydefault}{\gtdefault}
%\usepackage[haranoaji]{luatexja-preset}

% スライド番号の表示 %%%%%%%%%%%%%%%%%%%
\makeatletter
\setbeamertemplate{footline}{%
  \leavevmode%
  \hbox{%
  \begin{beamercolorbox}[wd=.5\paperwidth,ht=2.25ex,dp=1ex,right]{date in head/foot}%
    \insertframenumber{} \hspace*{2ex} % new version without total frames
  \end{beamercolorbox}}%
  \vskip0pt%
}
\makeatother
% スライド番号の表示 %%%%%%%%%%%%%%%%%%%
\usepackage{bookmark}
\IfFileExists{xurl.sty}{\usepackage{xurl}}{} % add URL line breaks if available
\urlstyle{same}
\hypersetup{
  hidelinks,
  pdfcreator={LaTeX via pandoc}}

\title{第3章: 多重回帰分析：推定}
\subtitle{Jeffrey Wooldridge (2016).\\
Introductory Econometrics: A Modern Approach\\
6rd ed.~Cengage Learning.}
\author{}
\date{\vspace{-2.5em}2026-01-23}

\begin{document}
\frame{\titlepage}

\section{準備}\label{ux6e96ux5099}

\begin{frame}[fragile]{必要なパッケージの読み込み}
\phantomsection\label{ux5fc5ux8981ux306aux30d1ux30c3ux30b1ux30fcux30b8ux306eux8aadux307fux8fbcux307f}
\begin{itemize}
\tightlist
\item
  \texttt{wooldridge}パッケージの読み込み
\end{itemize}

\begin{Shaded}
\begin{Highlighting}[]
\FunctionTok{library}\NormalTok{(wooldridge)}
\end{Highlighting}
\end{Shaded}
\end{frame}

\section{3-1
多重回帰分析の動機}\label{ux591aux91cdux56deux5e30ux5206ux6790ux306eux52d5ux6a5f}

\begin{frame}{単回帰分析の限界}
\phantomsection\label{ux5358ux56deux5e30ux5206ux6790ux306eux9650ux754c}
\begin{itemize}
\tightlist
\item
  単回帰モデル \(y = \beta_0 + \beta_1 x + u\) では、誤差項 \(u\) に
  \(y\) に影響を与える他のすべての要因が含まれている。
\item
  \(x\) と \(u\)
  が相関している場合（つまり、重要な変数が欠落している場合）、OLS推定量は\textbf{バイアス}を持つ（偏りがある）。
\item
  多重回帰分析を用いることで、他の要因を\textbf{明示的に}モデルに組み込み、制御することができる。
\end{itemize}
\end{frame}

\begin{frame}{多重回帰モデルの一般形}
\phantomsection\label{ux591aux91cdux56deux5e30ux30e2ux30c7ux30ebux306eux4e00ux822cux5f62}
\[ y = \beta_0 + \beta_1 x_1 + \beta_2 x_2 + \dots + \beta_k x_k + u \]

\begin{itemize}
\tightlist
\item
  \(\beta_0\): 切片
\item
  \(\beta_j\): 変数 \(x_j\) の係数 (\(j=1, \dots, k\))
\item
  \(u\): 誤差項
\end{itemize}
\end{frame}

\section{3-2
OLS推定量の解釈}\label{olsux63a8ux5b9aux91cfux306eux89e3ux91c8}

\begin{frame}{係数の解釈：セテリス・パリバス}
\phantomsection\label{ux4fc2ux6570ux306eux89e3ux91c8ux30bbux30c6ux30eaux30b9ux30d1ux30eaux30d0ux30b9}
\begin{itemize}
\tightlist
\item
  多重回帰モデルにおける係数 \(\beta_j\)
  は、\textbf{他のすべての独立変数を一定に保ったまま}（ceteris
  paribus）、\(x_j\) を1単位変化させたときの \(y\) の変化量を表す。
\end{itemize}

\[ \Delta y = \beta_j \Delta x_j \quad (\text{他の } x \text{ は固定}) \]

\begin{itemize}
\tightlist
\item
  これにより、実験データでなくても、他の要因の影響を統計的に制御した上での因果的効果を推定できる可能性がある。
\end{itemize}
\end{frame}

\begin{frame}{``部分''回帰としての解釈}
\phantomsection\label{ux90e8ux5206ux56deux5e30ux3068ux3057ux3066ux306eux89e3ux91c8}
\begin{itemize}
\tightlist
\item
  \(\beta_j\) の推定値は、単に \(y\) を \(x_j\) で回帰したものではない。
\item
  \(x_j\) と他の独立変数との相関関係を取り除いた上で、\(x_j\) が \(y\)
  に与える固有の影響を反映している。
\end{itemize}
\end{frame}

\section{3-3
Rによる多重回帰分析の実装}\label{rux306bux3088ux308bux591aux91cdux56deux5e30ux5206ux6790ux306eux5b9fux88c5}

\begin{frame}[fragile]{例: 賃金方程式の拡張 (\texttt{wage1})}
\phantomsection\label{ux4f8b-ux8cc3ux91d1ux65b9ux7a0bux5f0fux306eux62e1ux5f35-wage1}
\begin{itemize}
\item
  単回帰モデル： \[ \log(wage) = \beta_0 + \beta_1 educ + u \]
\item
  多重回帰モデル（経験年数 \texttt{exper} と勤続年数 \texttt{tenure}
  を追加）：
  \[ \log(wage) = \beta_0 + \beta_1 educ + \beta_2 exper + \beta_3 tenure + u \]
\end{itemize}
\end{frame}

\begin{frame}[fragile]{Rでの実行}
\phantomsection\label{rux3067ux306eux5b9fux884c}
\begin{Shaded}
\begin{Highlighting}[]
\FunctionTok{data}\NormalTok{(wage1)}

\CommentTok{\# 単回帰}
\NormalTok{model\_simple }\OtherTok{\textless{}{-}} \FunctionTok{lm}\NormalTok{(}\FunctionTok{log}\NormalTok{(wage) }\SpecialCharTok{\textasciitilde{}}\NormalTok{ educ, }\AttributeTok{data =}\NormalTok{ wage1)}

\CommentTok{\# 多重回帰}
\NormalTok{model\_multi }\OtherTok{\textless{}{-}} \FunctionTok{lm}\NormalTok{(}\FunctionTok{log}\NormalTok{(wage) }\SpecialCharTok{\textasciitilde{}}\NormalTok{ educ }\SpecialCharTok{+}\NormalTok{ exper }\SpecialCharTok{+}\NormalTok{ tenure, }
                  \AttributeTok{data =}\NormalTok{ wage1)}
\end{Highlighting}
\end{Shaded}
\end{frame}

\begin{frame}[fragile]{結果の比較}
\phantomsection\label{ux7d50ux679cux306eux6bd4ux8f03}
\begin{longtable}[]{@{}llrr@{}}
\caption{係数の比較}\tabularnewline
\toprule\noalign{}
& Variable & Simple & Multiple \\
\midrule\noalign{}
\endfirsthead
\toprule\noalign{}
& Variable & Simple & Multiple \\
\midrule\noalign{}
\endhead
(Intercept) & Intercept & 0.584 & 0.284 \\
educ & educ & 0.083 & 0.092 \\
exper & exper & NA & 0.004 \\
tenure & tenure & NA & 0.022 \\
\bottomrule\noalign{}
\end{longtable}

\begin{itemize}
\tightlist
\item
  教育年数（\texttt{educ}）の係数が、単回帰の \(0.083\) から多重回帰では
  \(0.092\) に変化した。
\item
  経験年数や勤続年数を一定に保つと、教育の収益率は少し高くなることが示唆される（ただし、この変化の大きさの評価は文脈による）。
\end{itemize}
\end{frame}

\section{\texorpdfstring{3-4 決定係数
\(R^2\)}{3-4 決定係数 R\^{}2}}\label{ux6c7aux5b9aux4fc2ux6570-r2}

\begin{frame}{多重回帰における \(R^2\)}
\phantomsection\label{ux591aux91cdux56deux5e30ux306bux304aux3051ux308b-r2}
\begin{itemize}
\tightlist
\item
  独立変数をモデルに追加すると、SSR（残差二乗和）は決して増加しないため、\(R^2\)（決定係数）は\textbf{決して減少しない}。
\end{itemize}

\[ R^2 = \frac{\text{SSE}}{\text{SST}} = 1 - \frac{\text{SSR}}{\text{SST}} \]

\begin{itemize}
\tightlist
\item
  そのため、変数を追加することの妥当性を判断するために \(R^2\)
  だけを用いるのは不適切である場合がある（第4章以降で修正 \(R^2\)
  などを扱う）。
\end{itemize}
\end{frame}

\begin{frame}[fragile]{決定係数の計算例}
\phantomsection\label{ux6c7aux5b9aux4fc2ux6570ux306eux8a08ux7b97ux4f8b}
\begin{Shaded}
\begin{Highlighting}[]
\FunctionTok{summary}\NormalTok{(model\_simple)}\SpecialCharTok{$}\NormalTok{r.squared}
\end{Highlighting}
\end{Shaded}

\begin{verbatim}
## [1] 0.1858065
\end{verbatim}

\begin{Shaded}
\begin{Highlighting}[]
\FunctionTok{summary}\NormalTok{(model\_multi)}\SpecialCharTok{$}\NormalTok{r.squared}
\end{Highlighting}
\end{Shaded}

\begin{verbatim}
## [1] 0.3160133
\end{verbatim}

\begin{itemize}
\tightlist
\item
  説明変数を追加することで、説明力が約18.6\%から30.6\%に向上した。
\end{itemize}
\end{frame}

\section{3-5 除外変数バイアス (Omitted Variable
Bias)}\label{ux9664ux5916ux5909ux6570ux30d0ux30a4ux30a2ux30b9-omitted-variable-bias}

\begin{frame}{除外変数バイアスとは}
\phantomsection\label{ux9664ux5916ux5909ux6570ux30d0ux30a4ux30a2ux30b9ux3068ux306f}
\begin{itemize}
\item
  本来モデルに含めるべき変数（真のモデルにおいて \(\beta \neq 0\)
  であり、かつ \(x\)
  と相関している変数）を除外して推定を行うと、OLS推定量はバイアスを持つ。
\item
  真のモデル: \(y = \beta_0 + \beta_1 x_1 + \beta_2 x_2 + u\)
\item
  推定モデル: \(y = \tilde{\beta}_0 + \tilde{\beta}_1 x_1 + v\)
\end{itemize}

このとき、\(\tilde{\beta}_1\) の期待値は以下のようになる。

\[ E(\tilde{\beta}_1) = \beta_1 + \beta_2 \frac{\text{Cov}(x_1, x_2)}{\text{Var}(x_1)} \]
\end{frame}

\begin{frame}{バイアスの方向}
\phantomsection\label{ux30d0ux30a4ux30a2ux30b9ux306eux65b9ux5411}
バイアスの符号（プラスかマイナスか）は、\(\beta_2\) の符号と、\(x_1\) と
\(x_2\) の相関（\(\text{Corr}(x_1, x_2)\)）によって決まる。

\begin{longtable}[]{@{}lll@{}}
\caption{除外変数バイアスの方向}\tabularnewline
\toprule\noalign{}
Corr.x1..x2. & beta2 & Bias \\
\midrule\noalign{}
\endfirsthead
\toprule\noalign{}
Corr.x1..x2. & beta2 & Bias \\
\midrule\noalign{}
\endhead
正 (+) & 正 (+) & 正 (+) \\
正 (+) & 負 (-) & 負 (-) \\
負 (-) & 正 (+) & 負 (-) \\
負 (-) & 負 (-) & 正 (+) \\
\bottomrule\noalign{}
\end{longtable}
\end{frame}

\section{3-6
OLS推定量の分散}\label{olsux63a8ux5b9aux91cfux306eux5206ux6563}

\begin{frame}{OLS推定量の分散}
\phantomsection\label{olsux63a8ux5b9aux91cfux306eux5206ux6563-1}
\begin{itemize}
\tightlist
\item
  独立変数が複数ある場合、\(\hat{\beta}_j\)
  の分散は以下の式で与えられる。
\end{itemize}

\[ \text{Var}(\hat{\beta}_j) = \frac{\sigma^2}{\text{SST}_j (1 - R_j^2)} \]

\begin{itemize}
\tightlist
\item
  \(\sigma^2\): 誤差項の分散
\item
  \(\text{SST}_j\): \(x_j\) の総変動 (\(\sum (x_{ij} - \bar{x}_j)^2\))
\item
  \(R_j^2\): \(x_j\)
  を\textbf{他のすべての独立変数}で回帰したときの決定係数
\end{itemize}
\end{frame}

\begin{frame}{多重共線性 (Multicollinearity)}
\phantomsection\label{ux591aux91cdux5171ux7ddaux6027-multicollinearity}
\begin{itemize}
\tightlist
\item
  \(R_j^2\) が1に近い場合（つまり、\(x_j\)
  が他の独立変数と強く相関している場合）、分母が0に近づき、\(\text{Var}(\hat{\beta}_j)\)
  が非常に大きくなる。
\item
  これを\textbf{多重共線性}の問題と呼ぶ。
\item
  分散が大きいと、推定精度が低くなり、t検定などで有意な結果が得られにくくなる。
\end{itemize}
\end{frame}

\begin{frame}{ガウス・マルコフの定理}
\phantomsection\label{ux30acux30a6ux30b9ux30deux30ebux30b3ux30d5ux306eux5b9aux7406}
\begin{itemize}
\tightlist
\item
  以下の5つの仮定（MLR.1 \textasciitilde{}
  MLR.5）の下で、OLS推定量は\textbf{最良線形不偏推定量}（BLUE: Best
  Linear Unbiased Estimator）となる。
\end{itemize}

\begin{enumerate}
\tightlist
\item
  \textbf{線形性}: パラメータについて線形である。
\item
  \textbf{ランダムサンプリング}: 無作為標本である。
\item
  \textbf{多重共線性なし}: 独立変数間に完全な線形関係がない。
\item
  \textbf{ゼロ条件付き平均}: \(E(u|x_1, \dots, x_k) = 0\)。
\item
  \textbf{均一分散}: \(\text{Var}(u|x_1, \dots, x_k) = \sigma^2\)。
\end{enumerate}
\end{frame}

\section{まとめ}\label{ux307eux3068ux3081}

\begin{frame}{まとめ}
\phantomsection\label{ux307eux3068ux3081-1}
\begin{itemize}
\tightlist
\item
  多重回帰分析は、他の要因を制御した上で変数の効果を推定できる。
\item
  OLS推定量は、セテリス・パリバス（他の条件一定）の効果として解釈する。
\item
  重要な変数が欠落していると、除外変数バイアスが生じる可能性がある。
\item
  独立変数間の相関が高いと、推定量の分散が大きくなる（多重共線性）。
\item
  ガウス・マルコフの定理により、特定の仮定の下でOLSは最も効率的な不偏推定量である。
\end{itemize}
\end{frame}

\end{document}
